\section*{A note on the Six Birds programme}
Six Birds' organizing picture motivates constructing upper computations over
lower-level objects \citep[\S1.1]{tsiokos2026foundations}. Whether
final-output training discovers such a construction is this programme's
empirical hypothesis, not a theorem about trained networks. Foundations I
separates stability, novelty and directionality certificates and keeps them
independent. Our five
achievements are different empirical measurements. Read in the theory's
terms, our central finding is that learning a strict extension induced the
lower theory only partially, and these separations help explain part of the
measured jagged competence; this is an empirical observation in one task, not a test of
the theory's formal certificate claims, and it does not establish that the
networks implement those layers internally.

\section*{Competing interests}
The author also developed Six Birds Theory, which motivates the task's
layered design and learning hypothesis; the empirical results stand
independently of the theory and bear on that application without validating
its formal claims. The author declares no other competing interests.

\section*{AI-assistance disclosure}
The author used AI tools in preparing this work and takes full responsibility
for its content.

\section*{Availability of materials}
The source package contains the manuscript, the claims ledger and evidence
manifest, the scripts that rebuild every number, figure and table, and the
released data files with their schema (\cref{app:compute}). The study code is
available at \url{https://github.com/ioannist/six-birds-ml}.
